% !TeX root = ../../Book4.tex
% 纲要标题：### 4.1 加性范畴
% 纲要位置：计划纲要.md，第 3188 行。
% * 分裂幂等元及其意义；Karoubi 包络及有限生成投射模

\section{加性范畴}
\label{b4:sec:0401}

模同态可以相加，也可以写成相对于有限直和的矩阵。将这些运算抽象出来以后，还须分别检查核、余核与幂等态射的像是否存在。

% 纲要内容（仅作写作计划，尚非教材正文）：
%
% * Hom 群结构与双积
% * 零对象、核与余核
%

\subsection{Hom 群结构与双积}
\label{b4:subsec:0401:01}

\begin{definition}\label{b4:def:additive-category}
设 \(\mathcal A\) 是局部小范畴。若每个 \(\operatorname{Hom}_{\mathcal A}(X,Y)\) 带有交换群结构，且态射复合对两侧加法均满足分配律，则称 \(\mathcal A\) 为\term[yujiaxing-fanchou]{预加性范畴}[preadditive category]。在预加性范畴 \(\mathcal A\) 中，给定对象 \(X,Y,B\) 及态射 \(i_X:X\to B\)、\(i_Y:Y\to B\)、\(p_X:B\to X\)、\(p_Y:B\to Y\)。若满足
\[
p_Xi_X=1_X,\quad p_Yi_Y=1_Y,\quad p_Xi_Y=p_Yi_X=0,
\qquad i_Xp_X+i_Yp_Y=1_B,
\]
则称 \(B\) 连同这些态射为 \(X,Y\) 的\term[shuangji]{双积}[biproduct]，记为 \(X\oplus Y\)。若预加性范畴还有零对象及所有有限双积，则称为\term[jiaxing-fanchou]{加性范畴}[additive category]。态射集中的零元称为零态射。复合的分配律保证 \(f0=0\) 与 \(0g=0\)。
\end{definition}

定义中的 \(0\) 不是另选的一组常值映射，而是态射群的零元。设零对象为 \(O\)，则 \(X\to O\to Y\) 的复合正是这个零态射：\(1_O=0\)，故经过 \(O\) 的任一复合都为零。没有零对象的预加性范畴仍有零态射。

\begin{proposition}\label{b4:prop:biproduct-calculus}
设 \(\mathcal A\) 为预加性范畴。\(\mathcal A\) 中任意双积连同其投影和包含同时满足积与余积的泛性质；反之，任意两个对象的积若存在，便可赋予唯一的与给定投影相容的双积结构。因此，有零对象及有限积的预加性范畴是加性范畴。有限双积之间的态射由其分量矩阵唯一决定，复合对应矩阵乘法。
\end{proposition}
\begin{proof}
给定 \(u:T\to X\)、\(v:T\to Y\)，唯一具有投影分量 \(u,v\) 的态射是 \(i_Xu+i_Yv\)：存在性由四个投影公式给出，唯一性由最后一个恒等式给出。给定 \(a:X\to T\)、\(b:Y\to T\)，余积要求的态射同样唯一地等于 \(ap_X+bp_Y\)。

反过来，若 \((P,p_X,p_Y)\) 是积，则由积的泛性质定义 \(i_X\) 的分量为 \((1_X,0)\)，\(i_Y\) 的分量为 \((0,1_Y)\)。\(i_Xp_X+i_Yp_Y\) 与 \(1_P\) 的两个投影相同，所以相等。这给出全部双积公式。空积是终对象；在预加性范畴中，若终对象 \(O\) 存在，则 \(1_O=0\)，任何 \(O\to X\) 都只能为零，所以 \(O\) 也是始对象。

对有限双积，态射 \(f:\bigoplus_jX_j\to\bigoplus_iY_i\) 的分量为 \(f_{ij}=p_if i_j\)，并有 \(f=\sum_{i,j}i_if_{ij}p_j\)。插入中间双积的恒等式 \(1=\sum_k i_kp_k\)，即得 \((gf)_{ij}=\sum_k g_{ik}f_{kj}\)。
\end{proof}

例如，任意含幺环 \(R\) 的左模范畴 \(R\text{-}\mathbf{Mod}\) 是加性范畴。有限直和里的投影和包含满足定义中的公式。有限维 \(k\)-向量空间及其线性映射也构成加性范畴，但非零向量空间组成的全子范畴没有零对象；不能因为其中仍能相加映射，就称它为加性范畴。

\subsection{零对象、核与余核}
\label{b4:subsec:0401:02}

有了零态射，“被某一映射送到零”便可表述为一个等式。为了定义核，不必先把对象看成元素的集合。

\begin{definition}\label{b4:def:kernel-cokernel}
设 \(\mathcal C\) 为有零态射的范畴，\(f:X\to Y\) 为其中的态射。给定态射 \(k:K\to X\)。若 \(fk=0\)，且对每个对象 \(T\in\mathcal C\) 及每个满足 \(ft=0\) 的态射 \(t:T\to X\)，存在唯一态射 \(u:T\to K\)，使 \(t=ku\)，则称 \(k\) 为 \(f\) 的\term[he]{核}[kernel]。给定态射 \(q:Y\to Q\)。若 \(qf=0\)，且对每个对象 \(T\in\mathcal C\) 及每个满足 \(sf=0\) 的态射 \(s:Y\to T\)，存在唯一态射 \(v:Q\to T\)，使 \(s=vq\)，则称 \(q\) 为 \(f\) 的\term[yuhe]{余核}[cokernel]。也将 \(K,Q\) 分别记为 \(\ker f,\operatorname{coker}f\)，但在复合公式中必须保留结构态射 \(k,q\)。
\end{definition}

核与余核的唯一分解分别由下图表示；虚线箭头由实线数据唯一确定：
\[
\begin{tikzcd}[column sep=2.5em,row sep=2.8em]
T\arrow[d,dashed,"u"']\arrow[dr,"t"]\arrow[drr,bend left=20,"0"]&&\\
K\arrow[r,"k"']&X\arrow[r,"f"']&Y
\end{tikzcd}
\qquad
\begin{tikzcd}[column sep=2.5em,row sep=2.8em]
X\arrow[r,"f"]\arrow[drr,bend right=20,"0"']&Y\arrow[r,"q"]\arrow[dr,"s"]&Q\arrow[d,dashed,"v"]\\
&&T
\end{tikzcd}
\]
核就是 \(f\) 与零态射的等化子，余核就是它们的余等化子。因此核态射是单态射，余核态射是满态射。这里的“单”“满”指态射的消去性质；在模范畴中才分别等价于集合映射的单射、满射。在加性范畴中，核的泛性质还说明：若核存在，则 \(f\) 为单态射当且仅当 \(\ker f=0\)。事实上，\(fu=fv\) 等价于 \(f(u-v)=0\)，这时 \(u-v\) 分解经过零对象。余核给出对偶结论。

\ProseParagraphBreak
在左 \(R\)-模中，这一定义恢复通常的子模 \(\{x:f(x)=0\}\) 及商模 \(Y/f(X)\)。在任意加性范畴中，定义却不保证这些对象存在。有限生成自由交换群构成的全子范畴就是一个需要谨慎的例子：它没有与交换群范畴相同的余核。\(\mathbb Z\xrightarrow{\times2}\mathbb Z\) 在这个子范畴中的余核反而是 \(\mathbb Z\to0\)，因为自由交换群没有非零的二阶元。\secref{b4:sec:0402}将说明这为何使它不能成为交换范畴。

\subsection{分裂幂等元及其意义}
\label{b4:subsec:0401:03}

若一个对象已经写成 \(X\oplus Y\)，复合 \(i_Xp_X\) 满足平方等于自身。反过来，一个这样的态射是否一定来自直和分解？这取决于其像能否作为范畴中的对象出现。

\begin{definition}\label{b4:def:split-idempotent}
设 \(\mathcal C\) 为范畴，\(e:X\to X\) 为其中的自态射。若 \(e^2=e\)，则称 \(e\) 为\term[mideng-taishe]{幂等态射}[idempotent morphism]。对幂等态射 \(e\)，若存在对象 \(Y\in\mathcal C\) 及态射 \(i:Y\to X\)、\(p:X\to Y\)，使 \(pi=1_Y\)、\(ip=e\)，则称它\term[fenlie-mideng]{分裂}[split]。
\end{definition}

\begin{proposition}\label{b4:prop:idempotent-kernel-splitting}
在加性范畴中，若幂等态射 \(e:X\to X\) 的 \(\ker e\) 和 \(\ker(1-e)\) 都存在，则
\[
X\cong\ker(1-e)\oplus\ker e,
\]
且 \(e\) 对应矩阵 \(\begin{pmatrix}1&0\\0&0\end{pmatrix}\)。特别地，每个具有全部核的加性范畴中，幂等态射都分裂。
\end{proposition}
\begin{proof}
记两个核态射为 \(i:Y\to X\)、\(j:Z\to X\)。由 \((1-e)e=0\) 和 \(e(1-e)=0\)，存在唯一 \(p:X\to Y\)、\(q:X\to Z\)，使 \(ip=e\)、\(jq=1-e\)。等式 \(ei=i\)、\(ej=0\) 与单态射 \(i,j\) 的消去性质给出
\[
pi=1_Y,\quad qj=1_Z,\quad pj=0,\quad qi=0.
\]
又有 \(ip+jq=1_X\)，故 \((i,j):Y\oplus Z\to X\) 与 \((p,q)^{\mathsf T}\) 互逆。最后，\(e i=i\)、\(e j=0\) 给出所说的矩阵。
\end{proof}

例如，在全部左 \(R\)-模中，幂等矩阵的像和核都是模，因而上述分解总存在。在有限秩自由左 \(R\)-模组成的范畴中，其像可能只是投射模，并非自由模。若 \(R=k\times k\)，则 \(R\to R\)、\(r\mapsto(1,0)r\) 是幂等态射，像为 \(k\times0\)；它不可能同构于 \(R^n\)，因为乘以 \((0,1)\) 后，前者为零，后者为 \(k^n\)。这个幂等态射不能在有限秩自由模范畴内分裂。这说明讨论分裂时必须指定所在范畴，而不能只检查矩阵满足 \(E^2=E\)。

\ProseParagraphBreak
缺少幂等态射的像时，可以把这个像作为一个新对象正式加入。原对象及其幂等态射已经说明了这个新对象应有的恒等态射，以及它与其他新对象之间哪些态射能够保留下来。这就是 Karoubi 包络的构造。\footnote{卡鲁比（Max Karoubi，1938-），法国数学家，主要研究代数 \(K\) 理论与拓扑 \(K\) 理论。}

\begin{definition}\label{b4:def:karoubi-envelope}
设 \(\mathcal C\) 为局部小范畴。若 \(\mathcal C\) 的每个幂等态射都分裂，则称 \(\mathcal C\) 为\term[mideng-wanbei]{幂等完备}[idempotent complete]的范畴。构造范畴 \(\operatorname{Kar}(\mathcal C)\)：对象为二元组 \((X,e)\)，其中 \(X\in\mathcal C\)、\(e:X\to X\) 且 \(e^2=e\)；从 \((X,e)\) 到 \((Y,d)\) 的态射为满足
\[
f=dfe
\]
的 \(\mathcal C\)-态射 \(f:X\to Y\)。复合沿用 \(\mathcal C\) 的复合，\((X,e)\) 的恒等态射为 \(e\)。这个范畴称为 \(\mathcal C\) 的\term[karoubi-baoluo]{Karoubi 包络}[Karoubi envelope]，也称幂等完备化。
\end{definition}

条件 \(f=dfe\) 等价于 \(df=f=fe\)。若再有 \(g=cgd\)，则 \(gf=c(gf)e\)，所以复合仍在指定的态射集中；上述两个等式也保证 \(e,d\) 分别在复合两端起恒等作用。由此定义确实给出范畴，而不只是对象与态射的列表。

\begin{proposition}\label{b4:prop:karoubi-splitting}
设 \(\mathcal C\) 为局部小范畴。函子
\[
J:\mathcal C\longrightarrow\operatorname{Kar}(\mathcal C),
\qquad X\longmapsto(X,1_X),\quad f\longmapsto f
\]
全忠实，且 \(\operatorname{Kar}(\mathcal C)\) 幂等完备。若 \(\mathcal C\) 加性，则其 Karoubi 包络也加性，且 \(J\) 为加性函子。
\end{proposition}
\begin{proof}
在 \((X,1_X)\) 与 \((Y,1_Y)\) 之间，态射条件不加任何限制，故 \(J\) 在 Hom 集上为恒等双射。

若 \(h\) 是 \((X,e)\) 上的幂等态射，则 \(h^2=h\) 且 \(eh=h=he\)。因此 \((X,h)\) 是一个对象，两条都由 \(h\) 给出的态射
\[
(X,h)\xrightarrow{h}(X,e)\xrightarrow{h}(X,h)
\]
的复合在 \((X,h)\) 上为恒等 \(h\)，反向复合则为 \((X,e)\) 上原来的 \(h\)。这就分裂了任意幂等态射。

若 \(\mathcal C\) 加性，条件 \(f=dfe\) 对加法封闭，故这些 Hom 集是原态射群的子群，复合仍双线性。零对象及双积分别由 \((0,1_0)\) 和 \((X\oplus Y,e\oplus d)\) 给出；把原双积的包含、投影限制到相应幂等分量，便满足双积公式。\(J\) 保持这些群运算与双积。
\end{proof}

\begin{theorem}\label{b4:thm:karoubi-universal-property}
设 \(\mathcal C\) 为小范畴，\(\mathcal D\) 为局部小、幂等完备的范畴，并为 \(\mathcal D\) 中每个幂等态射指定一个分裂，恒等态射的分裂取其自身。沿 \(J:\mathcal C\to\operatorname{Kar}(\mathcal C)\) 的限制给出等价
\[
J^*:\operatorname{Fun}(\operatorname{Kar}(\mathcal C),\mathcal D)
\longrightarrow\operatorname{Fun}(\mathcal C,\mathcal D).
\]
换言之，\(\mathcal C\) 上的函子可以延伸到其 Karoubi 包络，自然变换也唯一地延伸。
\end{theorem}
\begin{proof}
给定 \(F:\mathcal C\to\mathcal D\)。记 \(F(e)\) 的指定分裂为
\[
Y_{F,e}\xrightarrow{i_{F,e}}F(X)
\xrightarrow{p_{F,e}}Y_{F,e},
\qquad p_{F,e}i_{F,e}=1,\quad i_{F,e}p_{F,e}=F(e).
\]
定义 \(\overline F(X,e)=Y_{F,e}\)，对 \(f:(X,e)\to(Y,d)\) 定义
\[
\overline F(f)=p_{F,d}F(f)i_{F,e}.
\]
它把恒等 \(e\) 送到恒等。复合中插入 \(i_{F,d}p_{F,d}=F(d)\)，再用 \(gd=g\)、\(df=f\)，便得 \(\overline F(g)\overline F(f)=\overline F(gf)\)。恒等分裂的约定保证 \(\overline FJ=F\)。

若 \(\alpha:F\Rightarrow G\)，定义
\[
\overline\alpha_{(X,e)}
=p_{G,e}\alpha_Xi_{F,e}.
\]
\(\alpha_XF(e)=G(e)\alpha_X\) 保证它与分裂相容；在自然性方块中插入这两个幂等元，再用 \(\alpha_YF(f)=G(f)\alpha_X\)，就得到 \(\overline G(f)\overline\alpha_{(X,e)}=\overline\alpha_{(Y,d)}\overline F(f)\)。同一公式保持变换的恒等与复合，所以 \(F\mapsto\overline F\) 是限制函子的一个反方向函子。

还须说明为什么没有额外的自然变换。每个 \(E=(X,e)\) 都有收缩图
\[
E\xrightarrow{s_E}JX\xrightarrow{r_E}E,
\qquad r_Es_E=1_E,\quad s_Er_E=J(e),
\]
其中两条底层态射都为 \(e\)。给定 \(H,K:\operatorname{Kar}(\mathcal C)\to\mathcal D\) 及 \(\alpha:HJ\Rightarrow KJ\)，其延伸只能是
\[
\beta_E=K(r_E)\alpha_XH(s_E).
\]
这是因为自然性与 \(r_Es_E=1_E\) 已经迫使该等式成立。对 \(f:E\to E'\)，恒等式 \(f=r_{E'}J(f)s_E\) 及 \(\alpha\) 的自然性也表明这些 \(\beta_E\) 确实自然。因此限制全忠实。

最后，\(\overline{HJ}(E)\) 与 \(H(E)\) 都分裂 \(H(J(e))\)。若前一分裂记为 \(i_e,p_e\)，则
\[
H(r_E)i_e:\overline{HJ}(E)\to H(E),
\qquad p_eH(s_E):H(E)\to\overline{HJ}(E)
\]
互逆，且由刚证明的延伸唯一性组成自然同构。于是 \(F\mapsto\overline F\) 是 \(J^*\) 的拟逆。
\end{proof}

指定分裂是延伸函子所需的数据；例如模范畴可以直接取幂等同态的像。若目标范畴小，则通常选择公理足以同时选定全部分裂。加性范畴中的构造及其泛性质可对照 Bühler\cite[Section 6, Proposition 6.10, Corollary 6.11 and Example 6.12]{Buehler2010ExactCategories}。

\begin{proposition}\label{b4:prop:karoubi-free-projective}
设 \(R\) 为含幺环，\(\mathcal F_R\) 为以 \(R^n\)（\(n\ge0\)）为对象、以全部左 \(R\)-模同态为态射的范畴。映射
\[
\operatorname{Kar}(\mathcal F_R)\longrightarrow R\text{-}\mathbf{Mod},
\qquad (R^n,e)\longmapsto e(R^n)
\]
为全忠实函子，其对象像的同构类型恰为全部有限生成投射左 \(R\)-模。在为这些投射模指定有限自由模中的直和分解后，它给出与有限生成投射模范畴的等价。
\end{proposition}
\begin{proof}
幂等同态给出 \(R^n=e(R^n)\oplus(1-e)(R^n)\)，故其像是有限生成投射模。满足 \(f=dfe\) 的同态限制为 \(e(R^n)\to d(R^m)\)。反过来，给定这个像之间的同态 \(u\)，先用 \(e\) 投影、再作 \(u\)、最后包含入 \(R^m\)，便得到唯一满足 \(f=dfe\) 的 \(f\)。所以函子全忠实。

若 \(P\) 有限生成投射，则存在满同态 \(p:R^n\to P\)，投射性给出截面 \(i:P\to R^n\)。于是 \(e=ip\) 幂等，且 \(p:e(R^n)\to P\) 与 \(i\) 互逆。指定各 \(p,i\) 后，把 \(P\) 送到 \((R^n,ip)\)，把 \(u:P\to Q\) 送到 \(i_Qu p_P\)，就得到上述像函子的拟逆。
\end{proof}
